\chapter{Related Work}
\label{ch:related}

We organise related work around the thesis's two contributions. \Cref{sec:rwbench} reviews
the evaluation benchmarks that define the methodological standard the thesis must meet;
\cref{sec:rwagents} reviews the multi-agent trading frameworks whose architecture we adopt
and adapt; and \cref{sec:rwpos} states our position relative to both.

\section{Evaluation benchmarks: the methodological standard}
\label{sec:rwbench}

\paragraph{StockBench.} StockBench~\cite{stockbench2025} is the closest published setup to
ours: a contamination-free benchmark in which LLM agents make sequential buy/sell/hold
decisions in realistic multi-month equity environments with daily prices, fundamentals, and
news. We borrow its contamination-free framing and its return-plus-risk metric set, and its
equity-centric design makes it directly comparable to our primary US-equities track. Where we
extend it is on the reasoning axis: StockBench evaluates the \emph{actions} an agent takes,
whereas our central addition is to evaluate the \emph{chain-of-thought} behind them ---
faithfulness and grounding --- and to do so by regime.

\paragraph{LiveTradeBench.} LiveTradeBench~\cite{livetradebench2025} evaluates many LLMs in a
real-time window with a fixed lookback to prevent leakage. Two findings are load-bearing for
this thesis. First, general LLM leaderboard scores barely correlate with trading returns,
which justifies building a trading-specific evaluation rather than trusting generic
capability scores. Second, the large cross-regime generalisation gaps it reports motivate our
regime-segmented protocol directly. Our \mcp-as-time-machine design can be read as
generalising LiveTradeBench's fixed-lookback idea from a single global window to a per-tool,
per-timestamp filter that holds identically in backtest and live.

\paragraph{KellyBench.} KellyBench~\cite{kellybench2026} is the methodological guardrail. In a
long-horizon sequential-decision setting it documents that every frontier model loses money
on average and catalogues the precise failure modes a rigorous study must avoid: the
knowledge--action gap (a sizing rule coded but never invoked), non-stationarity, label
leakage from fitting on the full training set before the split, premature termination, and
single-seed unreliability. It also introduces a multi-point sophistication rubric. Our design
is a point-by-point response: faithfulness and grounding target the knowledge--action gap;
the \tnow\ filter and unadjusted prices target leakage; multi-seed evaluation with bootstrap
intervals targets single-seed unreliability; pre-registration of the faithfulness metric
targets post-hoc metric tuning; and our sophistication rubric is modelled on theirs. We cite
it heavily in the methodology and rigour sections.

\section{Multi-agent trading frameworks}
\label{sec:rwagents}

\paragraph{TradingAgents.} \textsc{TradingAgents}~\cite{tradingagents2024} is the reference
multi-agent ``firm'': fundamental, sentiment, and technical analysts; Bull and Bear
researchers in debate; a trader; and a risk team. We reuse its role decomposition and, in
particular, its Bull/Bear debate pattern --- the structure where a cyclic graph framework
genuinely pays off --- but adopt only the communication \emph{contract}, re-implementing it
over our \mcp\ tool layer rather than importing its code, so that our temporal filter, decision
schema, and evaluation plumbing remain in force. Our differentiator from it is exactly the two
contributions of this thesis: the \mcp\ tool layer with its structural look-ahead guarantee,
and the regime-aware reasoning evaluation.

\paragraph{Agent Market Arena.} AMA~\cite{ama2025} compares several agent designs across
several backbones in a live benchmark and reports that agent architecture drives behaviour
more than the model backbone does. This justifies two of our choices: comparing single-agent
against multi-agent architectures rather than sweeping many backbones, and treating the
multi-agent extension as a test of behavioural change rather than an assumed improvement.

\paragraph{ATLAS.} ATLAS~\cite{atlas2025} combines dynamic prompt optimisation with
multi-agent coordination and reports that its gains concentrate in \emph{volatile} regimes.
This directly supports our regime axis and the hypothesis that reasoning quality and strategy
must adapt to volatility; it is a useful contrast in that our contribution is to \emph{measure}
regime-dependent reasoning quality rather than to optimise prompts for it.

\paragraph{HedgeAgents and FinCon.} HedgeAgents~\cite{hedgeagents2025} (a central fund manager
coordinating asset-class experts via ``conferences'') and FinCon~\cite{fincon2024}
(hierarchical manager--analyst communication with risk control) inform the hierarchy of our
Portfolio-Manager extension, but we keep that extension deliberately narrow, consistent with
the evaluation-first priority.

\section{Positioning}
\label{sec:rwpos}

The reference architecture for multi-agent trading already exists; merely reproducing it would
add little. Our angle is orthogonal and matches the thesis topic precisely: a standardised
\mcp\ tool layer that makes look-ahead structurally impossible and an equivalent backtest/live
interface, combined with a regime-segmented evaluation of \emph{reasoning} quality
(faithfulness, grounding, sophistication) anchored to human judgement. Where prior benchmarks
evaluate what an agent \emph{does}, we additionally evaluate whether it does what it
\emph{says}, and we make the apparatus that measures this the primary object of study.
